\documentclass[12pt,a4paper]{article}

\usepackage[cm]{fullpage}
\usepackage{amsthm}
\usepackage{amsmath}
\usepackage{amsfonts}
\usepackage{amssymb}
\usepackage{xspace}
\usepackage[english]{babel}
\usepackage{fancyhdr}
\usepackage{titling}
\usepackage{framed}
\usepackage{stmaryrd}
\usepackage{mathpartir}
\renewcommand{\thesection}{Task \projnumber.\arabic{section}:}

%%%%%%%%%%%%%%%%%%%%%%%%%%%%%%%%%%%%%%%%%%%%%%%%%%%%%%%%
%         exercise-specific macros       %
%%%%%%%%%%%%%%%%%%%%%%%%%%%%%%%%%%%%%%%%%%%%%%%%%%%%%%%%

\newcommand{\xorshift}{\texttt{xorshift128+}}
\newcommand{\step}{\rightarrow}
\newcommand{\nstep}{\overset{*}{\rightarrow}}
\newcommand{\update}[3]{#1[#2 \mapsto #3]}
\newcommand{\run}[2]{\llbracket #1 \rrbracket #2}
\newcommand{\loweq}{\mathbin{=_L}}

\newcommand{\config}[1]{\ensuremath{\left\langle #1 \right\rangle}\xspace}
\newcommand{\ttt}[1]{\texttt{#1}}
\newcommand{\proves}{\vdash}
\newcommand{\sth}{\ensuremath{.\;}\xspace}
\newcommand\set[1]{\ensuremath{\{#1\}}}
\newcommand{\WHILE}[1]{\ensuremath{\mathsf{while}\ #1 \ \mathsf{do}}\xspace}
\newcommand{\IF}{\ensuremath{\mathsf{if}}\xspace}
\newcommand{\THEN}{\ensuremath{\mathsf{then}}\xspace}
\newcommand{\ELSE}{\ensuremath{\mathsf{else}}\xspace}

%%%%%%%%%%%%%%%%%%%%%%%%%%%%%%%%%%%%%%%%%%%%%%%%%%%%%%%%
% This part needs customization from you %
%%%%%%%%%%%%%%%%%%%%%%%%%%%%%%%%%%%%%%%%%%%%%%%%%%%%%%%%

% please enter your name and matriculation number here
\newcommand{\name}{Benjamin Deutsch}
\newcommand{\matriculation}{12215881}

%%%%%%%%%%%%%%%%%%%%%%%%%%%%%%%%%%%%%%%%%%%%%%%%%%%%%%%%
%           End of customization         %
%%%%%%%%%%%%%%%%%%%%%%%%%%%%%%%%%%%%%%%%%%%%%%%%%%%%%%%%

\newcommand{\projnumber}{1}
\newcommand{\Title}{Foundations: Z3 and Lean}
\setlength{\headheight}{15.2pt}
\setlength{\headsep}{20pt}
\setlength{\textheight}{680pt}
\pagestyle{fancy}
\fancyhf{}
\fancyhead[L]{Project \projnumber\ -- \Title}
\fancyhead[C]{}
\fancyhead[R]{\name{} (\matriculation)}
\renewcommand{\headrulewidth}{0.4pt}
\fancyfoot[C]{\thepage}


\begin{document}
\thispagestyle{empty}
\noindent\framebox[\linewidth]{%
 \begin{minipage}{\linewidth}%
 \hspace*{5pt} \textbf{Formal Methods for Security and Privacy (SS26)} \hfill Prof.~Matteo Maffei \hspace*{5pt}\\

 \begin{center}
  {\bf\Large Project \projnumber~-- \Title}
 \end{center}

 \vspace*{5pt}\hspace*{5pt}\name{} (\matriculation) \hfill TU Wien \hspace*{5pt}
 \end{minipage}%
 }
\vspace{0.5cm}
%%%%%%%%%%%%%%%%%%%%%%%%%%%%%%%%%%%%%%%%%%%%%%%%%%%%%%%%%%

% README
% remember to fill the macros \matriculation and \name

\subsection*{Project Task 4.2: Predicting the Next Random Numbers}
---

\subsection*{Project Task 4.3: Horn-Clause-Based Abstraction - DO-WHILE}
\begin{enumerate}
    \item We can encode the DO-WHILE loop using the following Horn clauses:
    \begin{align*}
        \{&S_{pc}(\bar v) \Rightarrow S_{pc'}(\bar v) \\
        &E_{pc'}(\bar v)\land \mathrm{EXPR}\neq 0 \Rightarrow S_{pc}(\bar v) \\
        &E_{pc'}(\bar v)\land \mathrm{EXPR} = 0 \Rightarrow E_{pc}(\bar v)\}
    \end{align*}

    These clauses can be translated into the following z3-python pseudo-code for $pc=1$ and $pc'=2$:
    \begin{verbatim}
    [
      ForAll([v_1,...,v_n],
        Implies( S1(v_1,...,v_n), S2(v_1,...,v_n) ),
      ForAll([v_1,...,v_n],
        Implies(
          And( E2(v_1,...,v_n), EXPR ),
          S1(v_1,...,v_n) ))
      ForAll([v_1,...,v_n],
        Implies(
          And( E2(v_1,...,v_n), Not(EXPR) ),
          E1(v_1,...,v_n) ))
    ]
    \end{verbatim}

    \item $S_{pc}(\bar v) \Rightarrow S_{pc'}(\bar v)$ makes sure that when we the loop is executed, the $\mathrm{CMD}$ is immediately executed. $E_{pc'}(\bar v)\land \mathrm{EXPR}\neq 0 \Rightarrow S_{pc}(\bar v)$ means that after execution of $\mathrm{CMD}$, given that $\mathrm{EXPR}$ evaluates to true, the loop is executed again. $E_{pc'}(\bar v)\land \mathrm{EXPR} = 0 \Rightarrow E_{pc}(\bar v)$ means that after execution of $\mathrm{CMD}$, given that $\mathrm{EXPR}$ evaluates to false, the loop is not executed again. This resembles the semantics of the DO-WHILE loop.

    \item Small-step semantics are rules that resemble a single step of computation. Given a program and a state, one evaluates the program execution step-by-step. In big-step semantics the evaluation of a program is described as a relation (often denoted as $\Downarrow$) between a program with its state and its final result. Using small-step semantics one can define the big-step semantics as the following: \[\langle \textrm{PROGRAM},\sigma\rangle \Downarrow \sigma' \Leftrightarrow \langle\textrm{PROGRAM},\sigma\rangle \overset{*}{\hookrightarrow} \langle\epsilon,\sigma'\rangle.\] 

    \item I would say that, since the premise $\langle \mathrm{CMD},\sigma\rangle \overset{*}{\hookrightarrow} \langle\epsilon,\sigma'\rangle$ can include multiple steps of computation, this rule is not a pure small-step rule. It is also not a pure big-step rule, since the conclusion relates programs with their states to programs with their states, like a small-step semantics rule.

    The rule has elements from both small-step and big-step semantics.
\end{enumerate}


\subsection*{Project Task 5.1: Proving Properties of xorshift}
---

\end{document}
